\section{Introduction}
\label{sec:intro}


Modern generative models~\citep{ldm,dit,sit, sd3}, trained on vast data using extensive computational resources, can be dramatically improved by aligning their internal features with those of a frozen image encoder, for example, the 86M parameter model \emph{DINO} ~\citep{repa}. This encoder was trained not to generate, but to \emph{discriminate}, i.e., to cluster images by semantic similarity. Its effectiveness for generative modeling exposes a gap: flow models do not learn strong representations on their own, although they help generation. External alignment offers a practical remedy: borrowing representations from a model that \emph{did} learn them. 

However, this approach has fundamental limitations. First, external alignment fails to uphold expected scaling laws, with stronger encoders often exhibiting diminished or even negative returns~\citep{irepa} (Sec.~\ref{sec:motivation}). Moreover, as we demonstrate in this work (see Sec.~\ref{sec:experiments}), scaling the generative model does not yield proportional improvements when relying on external alignment. Second, these methods fail to generalize across modalities: for video and audio generation, we find that alignment with most external encoders actually harms performance (Sec.~\ref{sec:experiments}), making external alignment less suitable for multi-modal models that must handle diverse data distributions within a single framework. Finally, it is difficult to anticipate which encoder will be effective for a given task. For example, aligning text-to-image models with SigLIP 2~\citep{siglip2} performs worse than DINOv2~\citep{dinov2} (Sec.~\ref{sec:experiments}), despite the former being explicitly trained with text supervision and multi-aspect ratio support, properties seemingly better suited for the task.

To avoid the use of external representations, existing approaches opt to rely on the model's natively learned features and the semantic asymmetry between different layers~\citep{sra,haghighi2025layersyncselfaligningintermediatelayers}. However, such formulations remain limited by the semantics naturally learned by the flow objective, and lag behind external alignment. 

In contrast to both approaches above, we propose to directly integrate a self-supervised framework into flow matching to actively strengthen representations beyond those learned by the generative objective alone.
To this end, we propose \emph{Dual-Timestep Scheduling}, which applies two distinct noise levels %sampled from the model's noise distribution 
to different subsets of input tokens, creating an information asymmetry in which some tokens are more heavily corrupted than others. 
We perform two forward passes: one with the mixed, heterogeneously-noised input, and one with a cleaner input where all tokens are noised at the lower of the two levels.
The self-supervised objective is to predict, from the mixed input, the representations the model produces given the cleaner input. Combined with the standard flow loss on the heterogeneously-noised input, the model thus learns both dense, flow-based reconstruction and semantic feature prediction within a unified framework.

Since our formulation operates purely on the model's internal representations without relying on external encoders, it naturally extends to both single-modality and joint multi-modal training. Fig.~\ref{fig:teaser} shows qualitative examples from our jointly-trained image, video, and audio model. Compared to standard flow matching, our method yields improvements in structural coherence, particularly for challenging structures like faces and hands, as well as text rendering accuracy and temporal consistency in video. Moreover, our method is agnostic to autoencoder choice: we demonstrate consistent improvements across SD~\cite{ldm}, FLUX.2~\cite{bfl2025representation}, Wan2.2~\cite{wan2025}, Songbloom~\cite{yang2025songbloom}, and representation autoencoders~\cite{rae} (Sec.~\ref{sec:experiments}).

We evaluate on image, video, audio, and multi-modal generation, and show that our framework outperforms leading external alignment methods in each setting. These results suggest that joint optimization of generation and representations offers a robust, scalable, and general path forward.